\chapter{Half Adders}
\label{chap:half-adders}

\chapterintro{A half adder is one of the simplest arithmetic circuits. It adds two binary digits and produces a sum bit and a carry bit.}

\learningobjectives{
\begin{itemize}
    \item Explain the purpose of a half adder.
    \item Derive the half adder truth table.
    \item Understand why XOR produces the sum bit.
    \item Understand why AND produces the carry bit.
\end{itemize}}

\section{The Problem}

Binary addition begins with the simplest case: adding two single bits.

There are four possibilities:

\begin{center}
\begin{tabular}{cccc}
\toprule
A & B & Sum & Carry \\
\midrule
0 & 0 & 0 & 0 \\
0 & 1 & 1 & 0 \\
1 & 0 & 1 & 0 \\
1 & 1 & 0 & 1 \\
\bottomrule
\end{tabular}
\end{center}

When \(1+1\), the result is \(10_2\). The rightmost bit is the sum and the leftmost bit is the carry.

\section{Deriving the Circuit}

The sum bit is \(1\) when the inputs are different, so:

\[
\text{Sum}=A\operatorname{XOR}B
\]

The carry bit is \(1\) only when both inputs are \(1\), so:

\[
\text{Carry}=A\operatorname{AND}B
\]

\section{Python Implementation}

\begin{lstlisting}
def half_adder(a: int, b: int) -> tuple[int, int]:
    if a not in (0, 1) or b not in (0, 1):
        raise ValueError("inputs must be 0 or 1")
    total = a + b
    return total % 2, total // 2
\end{lstlisting}

\section{Worked Example}

For \(A=1\) and \(B=1\):

\[
\text{Sum}=1\operatorname{XOR}1=0
\]

\[
\text{Carry}=1\operatorname{AND}1=1
\]

Therefore the result is \(10_2\).

\section{Exercises}

\begin{exercise}
Write the half adder output for \(A=0\), \(B=1\).
\end{exercise}

\begin{solution}
The sum is \(1\) and the carry is \(0\). The result is \(01_2\).
\end{solution}

\section{Portfolio Milestone}

Extend your logic simulator with a \code{half\_adder} function and tests for all four input combinations.

\section{Summary}

A half adder adds two bits. It uses XOR for the sum and AND for the carry.
